\chapter{Hypertension (High Blood Pressure)}
\index{Hypertension (High Blood Pressure)}

\section*{Definition and Overview}
Hypertension, or high blood pressure, is a condition in which the pressure of blood against the walls of the arteries is persistently too high. It is measured as two numbers: the systolic pressure, the peak pressure when the heart contracts, and the diastolic pressure, the pressure when the heart relaxes. Hypertension is usually defined as a resting blood pressure of 140/90 mmHg or higher, measured repeatedly. Because it usually causes no symptoms, many people do not know they have it, yet it is a major risk factor for heart attack, stroke, kidney disease, and dementia. Detection is simple, and treatment with lifestyle change and medicines greatly reduces these risks. Hypertension is sometimes called the silent killer for this reason.

\section*{Epidemiology}
Hypertension is extremely common worldwide and becomes more frequent with age. Many people are unaware they have it, and diagnosis and treatment rates vary between countries and communities. Risk is associated with family history, excess body weight, high salt intake, physical inactivity, tobacco, alcohol, sleep apnoea, kidney disease, and some medicines.

\section*{Aetiology and Pathophysiology}
Blood pressure is determined by the amount of blood the heart pumps and the resistance of the arteries, which is controlled by the diameter of the small arteries and by the balance of hormones and the nervous system. In most people no single cause is found; this is called primary or essential hypertension, and it develops gradually with age under the influence of genetics, excess weight, high salt intake, alcohol, inactivity, and stress. A minority have secondary hypertension from an identifiable cause, such as kidney disease, narrowing of the renal arteries, hormonal conditions such as primary aldosteronism or thyroid disease, sleep apnoea, or medicines including the contraceptive pill and some painkillers. Sustained high pressure damages the walls of arteries throughout the body, promoting atherosclerosis, and strains the heart, kidneys, eyes, and brain.

\section*{Clinical Features}
\begin{itemize}
    \item Hypertension usually causes no symptoms, even at very high levels
    \item It is often discovered during a routine check, health screening, or consultation for another reason
    \item Occasional non-specific symptoms such as headache or nosebleeds are not reliable indicators
    \item Very high blood pressure (a hypertensive crisis) can cause severe headache, visual disturbance, chest pain, breathlessness, or confusion
    \item Complications may be the first presentation: heart attack, stroke, kidney failure, or heart failure
    \item Physical findings can include a raised blood pressure reading and signs of target organ damage
\end{itemize}

\section*{Differential Diagnosis}
A single high reading does not mean hypertension, and the cause of consistently high readings should be considered.
\begin{itemize}
    \item \textbf{White-coat hypertension:} high readings in the clinic but normal at home, due to anxiety
    \item \textbf{Masked hypertension:} normal clinic readings but high at home
    \item \textbf{Measurement error:} incorrect cuff size, technique, or recent caffeine or exercise
    \item \textbf{Secondary hypertension:} kidney disease, hormonal conditions, sleep apnoea, and medicines
    \item \textbf{Transient rises:} pain, anxiety, or acute illness
    \item \textbf{Postural variation:} differences between lying, sitting, and standing
\end{itemize}

\section*{When to Seek Medical Care}
Adults should have their blood pressure checked regularly, at least every two years from age 18 and more often with risk factors. Anyone with a high reading should have it confirmed with repeat measurements, including home monitoring. People with known hypertension should have regular reviews of their blood pressure, medicines, and risk factors.

\textbf{Contact your local emergency services for:}
\begin{itemize}
    \item A very high blood pressure reading with severe headache, blurred vision, chest pain, or breathlessness
    \item Sudden weakness, numbness, facial droop, or difficulty speaking, which may indicate a stroke
    \item Chest pain or pressure, which may indicate a heart attack
    \item Collapse, confusion, or a seizure
\end{itemize}

\section*{Diagnosis and Investigations}
\begin{itemize}
    \item \textbf{Blood pressure measurement:} repeated readings with a validated cuff, on both arms, and at more than one visit
    \item \textbf{Home blood pressure monitoring:} readings over several days to confirm the diagnosis
    \item \textbf{Ambulatory monitoring:} 24-hour monitoring, useful for white-coat and masked hypertension
    \item \textbf{Urine and blood tests:} kidney function, electrolytes, glucose, and lipids
    \item \textbf{Electrocardiogram:} assessment of the heart for strain or enlargement
    \item \textbf{Cardiovascular risk assessment:} estimating the overall risk of heart attack and stroke
    \item \textbf{Investigations for secondary causes:} hormonal tests, ultrasound, or imaging where suspected
\end{itemize}

\section*{Management}
\begin{itemize}
    \item \textbf{Lifestyle measures:} the foundation of treatment, including healthy eating, reduced salt, weight loss, activity, and limiting alcohol
    \item \textbf{Medicines:} a range of classes, including ACE inhibitors, angiotensin receptor blockers, calcium channel blockers, and diuretics
    \item \textbf{Combination therapy:} most people need more than one medicine to reach target
    \item \textbf{Regular review:} monitoring of blood pressure and adjustment of treatment
    \item \textbf{Treatment of other risk factors:} cholesterol-lowering medicines and aspirin where indicated
    \item \textbf{Management of secondary causes:} treatment of kidney, hormonal, or sleep conditions
    \item \textbf{Blood pressure targets:} individualised, usually below 140/90 mmHg, and lower in some groups
    \item \textbf{Support and self-management:} home monitoring, written action plans, and pharmacist involvement
\end{itemize}

\section*{Complications}
\begin{itemize}
    \item Heart attack, heart failure, and angina
    \item Stroke and transient ischaemic attack
    \item Chronic kidney disease and kidney failure
    \item Damage to the eyes with visual loss
    \item Peripheral arterial disease
    \item Aortic aneurysm and dissection
    \item Cognitive decline and dementia
    \item Hypertensive crisis and its organ damage
    \item Side effects of blood pressure medicines, affecting adherence
\end{itemize}

\section*{Prognosis and Outcomes}
Treating hypertension substantially reduces the risk of stroke, heart attack, heart failure, and kidney disease, and the benefits are greatest in people with higher risk. Each reduction of 10 mmHg in systolic pressure cuts cardiovascular risk by around 20 per cent. Most people can control their blood pressure with a combination of lifestyle change and medicines, and most require two or more medicines. Treatment is lifelong, but it is well tolerated, and with good control, people with hypertension can expect a near-normal life expectancy.

\section*{Prevention and Self-Care}
\begin{itemize}
    \item Have your blood pressure checked regularly
    \item Eat a diet rich in fruit, vegetables, and whole grains, and low in salt and processed food
    \item Reduce salt to less than 5 grams (one teaspoon) per day
    \item Maintain a healthy weight and be physically active most days
    \item Limit alcohol and avoid smoking
    \item Manage stress and get enough sleep
    \item Take medicines exactly as prescribed, even when you feel well
    \item Monitor your blood pressure at home if advised
    \item Attend regular reviews and tell your doctor about any new medicines
\end{itemize}

\section*{Sources and Further Reading}
The following reputable sources provide further information for healthcare students and patients seeking reliable, current advice about high blood pressure.
\begin{itemize}
    \item World Health Organization. \textit{Hypertension.} https://www.who.int/health-topics/hypertension
    \item NHS. \textit{High blood pressure (hypertension).} https://www.nhs.uk/conditions/high-blood-pressure-hypertension/
\end{itemize}
